% F1 — Architecture overview (TikZ, paper-quality vector figure)
% Three-tier corpus + merge engine + LLM consumer.
% Color-coded by tier; dashed boundaries group conceptual regions.
% Style follows Gao et al. (Amazon Science 2022, RAMKG) Figure 1.
%
% Used in §3 Architecture as Figure 1.
% Include via: % F1 — Architecture overview (TikZ, paper-quality vector figure)
% Three-tier corpus + merge engine + LLM consumer.
% Color-coded by tier; dashed boundaries group conceptual regions.
% Style follows Gao et al. (Amazon Science 2022, RAMKG) Figure 1.
%
% Used in §3 Architecture as Figure 1.
% Include via: % F1 — Architecture overview (TikZ, paper-quality vector figure)
% Three-tier corpus + merge engine + LLM consumer.
% Color-coded by tier; dashed boundaries group conceptual regions.
% Style follows Gao et al. (Amazon Science 2022, RAMKG) Figure 1.
%
% Used in §3 Architecture as Figure 1.
% Include via: % F1 — Architecture overview (TikZ, paper-quality vector figure)
% Three-tier corpus + merge engine + LLM consumer.
% Color-coded by tier; dashed boundaries group conceptual regions.
% Style follows Gao et al. (Amazon Science 2022, RAMKG) Figure 1.
%
% Used in §3 Architecture as Figure 1.
% Include via: \input{figures/F1-architecture.tex}

\begin{figure}[t]
\centering
\resizebox{0.95\textwidth}{!}{%
\begin{tikzpicture}[
  font=\sffamily\small,
  >=stealth,
  node distance=0.6cm and 0.9cm,
  every node/.style={align=center},
  tier/.style={rectangle, rounded corners=2pt, draw, thick,
               minimum width=4.2cm, minimum height=1.3cm, inner sep=4pt},
  curated/.style={tier, fill=blue!10, draw=blue!60!black},
  super/.style={tier, fill=orange!18, draw=orange!70!black},
  digest/.style={tier, fill=gray!10, draw=gray!60!black},
  engine/.style={rectangle, draw, thick, fill=green!12, draw=green!50!black,
                 minimum width=3.4cm, minimum height=2.2cm, rounded corners=4pt},
  consumer/.style={cylinder, draw, thick, fill=violet!10, draw=violet!60!black,
                   shape border rotate=90, aspect=0.3,
                   minimum width=2.6cm, minimum height=1.3cm},
  group/.style={draw, dashed, gray, rounded corners=4pt, inner sep=6pt},
  flow/.style={->, thick, draw=black!70},
  legend/.style={font=\sffamily\scriptsize, gray!60!black}
]

% --- Tier nodes ---
\node[curated] (curated) {\textbf{Tier 1 — Curated memory}\\\scriptsize \texttt{*.md} (canonical)};
\node[super, below=of curated] (super) {\textbf{Supersession sidecars}\\\scriptsize \texttt{*.md.supersessions.jsonl}\\\scriptsize (append-only)};
\node[digest, below=of super] (digest) {\textbf{Tier 2 — Journal digests}\\\scriptsize transcript-distilled summaries};

% --- Dashed boundary grouping the curated + supersession ---
\node[group, fit=(curated)(super), label={[legend,anchor=south west,xshift=0pt,yshift=-2pt]north west:Operator-authored corpus}] (curatedgroup) {};

% --- Merge engine ---
\node[engine, right=2.2cm of super.east, anchor=west] (engine) {%
  \textbf{Retrieval Engine}\\[2pt]
  \scriptsize SQLite FTS5 + BM25\\
  \scriptsize \textsl{accretive merge}\\
  \scriptsize precedence: super $\prec$ canonical
};

% --- LLM consumer ---
\node[consumer, right=1.4cm of engine.east, anchor=west] (llm) {%
  \textbf{LLM}\\[1pt]\scriptsize agent / CDSS
};

% --- File-system watcher (small, top-right of engine) ---
\node[draw, thick, fill=yellow!18, draw=yellow!60!black,
      rectangle, rounded corners=2pt, minimum width=2cm, minimum height=0.7cm,
      above=0.5cm of engine.north, anchor=south,
      font=\sffamily\scriptsize] (watcher) {fs-watcher (1\,s)};

% --- Arrows ---
\draw[flow] (curated.east) -- ++(0.5cm,0) |- (engine.west |- curated);
\draw[flow] (super.east)   -- (engine.west |- super);
\draw[flow] (digest.east)  -- ++(0.5cm,0) |- (engine.west |- digest);

\draw[flow] (engine.east) -- node[above,font=\sffamily\scriptsize]{merged result\\(super first)} (llm.west);

\draw[flow, dashed, gray] (watcher.south) -- (engine.north);
\draw[flow, dashed, gray] (watcher.west) -| ++(-1cm,-0.3cm) -| (curated.north);

% --- Legend ---
\node[legend, below=0.4cm of digest.south, anchor=north, align=left, text width=12cm] {%
  \textbf{Precedence invariant:} for any retrieved file with $k \geq 1$ supersession records,
  the merged output places all $k$ records before the canonical body —
  ensuring sequential readers (LLMs or humans) encounter the correction first.\\
  \textit{Theorem~\ref{thm:precedence}: empirically verified at $100\%$ over $N{=}1000$ randomized queries.}
};

\end{tikzpicture}%
}
\caption{Architecture of memory-oracle. The retrieval engine indexes three corpus
tiers (curated files, supersession sidecars, journal digests) and merges
supersession records ahead of canonical bodies at read time. A filesystem
watcher absorbs new writes within $\sim$1\,s, enabling a self-extending
corpus. The LLM consumer (clinical decision support, coding agent, trading
strategy assessor) receives a single merged document in which corrections
structurally precede the assertions they supersede.}
\label{fig:architecture}
\end{figure}


\begin{figure}[t]
\centering
\resizebox{0.95\textwidth}{!}{%
\begin{tikzpicture}[
  font=\sffamily\small,
  >=stealth,
  node distance=0.6cm and 0.9cm,
  every node/.style={align=center},
  tier/.style={rectangle, rounded corners=2pt, draw, thick,
               minimum width=4.2cm, minimum height=1.3cm, inner sep=4pt},
  curated/.style={tier, fill=blue!10, draw=blue!60!black},
  super/.style={tier, fill=orange!18, draw=orange!70!black},
  digest/.style={tier, fill=gray!10, draw=gray!60!black},
  engine/.style={rectangle, draw, thick, fill=green!12, draw=green!50!black,
                 minimum width=3.4cm, minimum height=2.2cm, rounded corners=4pt},
  consumer/.style={cylinder, draw, thick, fill=violet!10, draw=violet!60!black,
                   shape border rotate=90, aspect=0.3,
                   minimum width=2.6cm, minimum height=1.3cm},
  group/.style={draw, dashed, gray, rounded corners=4pt, inner sep=6pt},
  flow/.style={->, thick, draw=black!70},
  legend/.style={font=\sffamily\scriptsize, gray!60!black}
]

% --- Tier nodes ---
\node[curated] (curated) {\textbf{Tier 1 — Curated memory}\\\scriptsize \texttt{*.md} (canonical)};
\node[super, below=of curated] (super) {\textbf{Supersession sidecars}\\\scriptsize \texttt{*.md.supersessions.jsonl}\\\scriptsize (append-only)};
\node[digest, below=of super] (digest) {\textbf{Tier 2 — Journal digests}\\\scriptsize transcript-distilled summaries};

% --- Dashed boundary grouping the curated + supersession ---
\node[group, fit=(curated)(super), label={[legend,anchor=south west,xshift=0pt,yshift=-2pt]north west:Operator-authored corpus}] (curatedgroup) {};

% --- Merge engine ---
\node[engine, right=2.2cm of super.east, anchor=west] (engine) {%
  \textbf{Retrieval Engine}\\[2pt]
  \scriptsize SQLite FTS5 + BM25\\
  \scriptsize \textsl{accretive merge}\\
  \scriptsize precedence: super $\prec$ canonical
};

% --- LLM consumer ---
\node[consumer, right=1.4cm of engine.east, anchor=west] (llm) {%
  \textbf{LLM}\\[1pt]\scriptsize agent / CDSS
};

% --- File-system watcher (small, top-right of engine) ---
\node[draw, thick, fill=yellow!18, draw=yellow!60!black,
      rectangle, rounded corners=2pt, minimum width=2cm, minimum height=0.7cm,
      above=0.5cm of engine.north, anchor=south,
      font=\sffamily\scriptsize] (watcher) {fs-watcher (1\,s)};

% --- Arrows ---
\draw[flow] (curated.east) -- ++(0.5cm,0) |- (engine.west |- curated);
\draw[flow] (super.east)   -- (engine.west |- super);
\draw[flow] (digest.east)  -- ++(0.5cm,0) |- (engine.west |- digest);

\draw[flow] (engine.east) -- node[above,font=\sffamily\scriptsize]{merged result\\(super first)} (llm.west);

\draw[flow, dashed, gray] (watcher.south) -- (engine.north);
\draw[flow, dashed, gray] (watcher.west) -| ++(-1cm,-0.3cm) -| (curated.north);

% --- Legend ---
\node[legend, below=0.4cm of digest.south, anchor=north, align=left, text width=12cm] {%
  \textbf{Precedence invariant:} for any retrieved file with $k \geq 1$ supersession records,
  the merged output places all $k$ records before the canonical body —
  ensuring sequential readers (LLMs or humans) encounter the correction first.\\
  \textit{Theorem~\ref{thm:precedence}: empirically verified at $100\%$ over $N{=}1000$ randomized queries.}
};

\end{tikzpicture}%
}
\caption{Architecture of memory-oracle. The retrieval engine indexes three corpus
tiers (curated files, supersession sidecars, journal digests) and merges
supersession records ahead of canonical bodies at read time. A filesystem
watcher absorbs new writes within $\sim$1\,s, enabling a self-extending
corpus. The LLM consumer (clinical decision support, coding agent, trading
strategy assessor) receives a single merged document in which corrections
structurally precede the assertions they supersede.}
\label{fig:architecture}
\end{figure}


\begin{figure}[t]
\centering
\resizebox{0.95\textwidth}{!}{%
\begin{tikzpicture}[
  font=\sffamily\small,
  >=stealth,
  node distance=0.6cm and 0.9cm,
  every node/.style={align=center},
  tier/.style={rectangle, rounded corners=2pt, draw, thick,
               minimum width=4.2cm, minimum height=1.3cm, inner sep=4pt},
  curated/.style={tier, fill=blue!10, draw=blue!60!black},
  super/.style={tier, fill=orange!18, draw=orange!70!black},
  digest/.style={tier, fill=gray!10, draw=gray!60!black},
  engine/.style={rectangle, draw, thick, fill=green!12, draw=green!50!black,
                 minimum width=3.4cm, minimum height=2.2cm, rounded corners=4pt},
  consumer/.style={cylinder, draw, thick, fill=violet!10, draw=violet!60!black,
                   shape border rotate=90, aspect=0.3,
                   minimum width=2.6cm, minimum height=1.3cm},
  group/.style={draw, dashed, gray, rounded corners=4pt, inner sep=6pt},
  flow/.style={->, thick, draw=black!70},
  legend/.style={font=\sffamily\scriptsize, gray!60!black}
]

% --- Tier nodes ---
\node[curated] (curated) {\textbf{Tier 1 — Curated memory}\\\scriptsize \texttt{*.md} (canonical)};
\node[super, below=of curated] (super) {\textbf{Supersession sidecars}\\\scriptsize \texttt{*.md.supersessions.jsonl}\\\scriptsize (append-only)};
\node[digest, below=of super] (digest) {\textbf{Tier 2 — Journal digests}\\\scriptsize transcript-distilled summaries};

% --- Dashed boundary grouping the curated + supersession ---
\node[group, fit=(curated)(super), label={[legend,anchor=south west,xshift=0pt,yshift=-2pt]north west:Operator-authored corpus}] (curatedgroup) {};

% --- Merge engine ---
\node[engine, right=2.2cm of super.east, anchor=west] (engine) {%
  \textbf{Retrieval Engine}\\[2pt]
  \scriptsize SQLite FTS5 + BM25\\
  \scriptsize \textsl{accretive merge}\\
  \scriptsize precedence: super $\prec$ canonical
};

% --- LLM consumer ---
\node[consumer, right=1.4cm of engine.east, anchor=west] (llm) {%
  \textbf{LLM}\\[1pt]\scriptsize agent / CDSS
};

% --- File-system watcher (small, top-right of engine) ---
\node[draw, thick, fill=yellow!18, draw=yellow!60!black,
      rectangle, rounded corners=2pt, minimum width=2cm, minimum height=0.7cm,
      above=0.5cm of engine.north, anchor=south,
      font=\sffamily\scriptsize] (watcher) {fs-watcher (1\,s)};

% --- Arrows ---
\draw[flow] (curated.east) -- ++(0.5cm,0) |- (engine.west |- curated);
\draw[flow] (super.east)   -- (engine.west |- super);
\draw[flow] (digest.east)  -- ++(0.5cm,0) |- (engine.west |- digest);

\draw[flow] (engine.east) -- node[above,font=\sffamily\scriptsize]{merged result\\(super first)} (llm.west);

\draw[flow, dashed, gray] (watcher.south) -- (engine.north);
\draw[flow, dashed, gray] (watcher.west) -| ++(-1cm,-0.3cm) -| (curated.north);

% --- Legend ---
\node[legend, below=0.4cm of digest.south, anchor=north, align=left, text width=12cm] {%
  \textbf{Precedence invariant:} for any retrieved file with $k \geq 1$ supersession records,
  the merged output places all $k$ records before the canonical body —
  ensuring sequential readers (LLMs or humans) encounter the correction first.\\
  \textit{Theorem~\ref{thm:precedence}: empirically verified at $100\%$ over $N{=}1000$ randomized queries.}
};

\end{tikzpicture}%
}
\caption{Architecture of memory-oracle. The retrieval engine indexes three corpus
tiers (curated files, supersession sidecars, journal digests) and merges
supersession records ahead of canonical bodies at read time. A filesystem
watcher absorbs new writes within $\sim$1\,s, enabling a self-extending
corpus. The LLM consumer (clinical decision support, coding agent, trading
strategy assessor) receives a single merged document in which corrections
structurally precede the assertions they supersede.}
\label{fig:architecture}
\end{figure}


\begin{figure}[t]
\centering
\resizebox{0.95\textwidth}{!}{%
\begin{tikzpicture}[
  font=\sffamily\small,
  >=stealth,
  node distance=0.6cm and 0.9cm,
  every node/.style={align=center},
  tier/.style={rectangle, rounded corners=2pt, draw, thick,
               minimum width=4.2cm, minimum height=1.3cm, inner sep=4pt},
  curated/.style={tier, fill=blue!10, draw=blue!60!black},
  super/.style={tier, fill=orange!18, draw=orange!70!black},
  digest/.style={tier, fill=gray!10, draw=gray!60!black},
  engine/.style={rectangle, draw, thick, fill=green!12, draw=green!50!black,
                 minimum width=3.4cm, minimum height=2.2cm, rounded corners=4pt},
  consumer/.style={cylinder, draw, thick, fill=violet!10, draw=violet!60!black,
                   shape border rotate=90, aspect=0.3,
                   minimum width=2.6cm, minimum height=1.3cm},
  group/.style={draw, dashed, gray, rounded corners=4pt, inner sep=6pt},
  flow/.style={->, thick, draw=black!70},
  legend/.style={font=\sffamily\scriptsize, gray!60!black}
]

% --- Tier nodes ---
\node[curated] (curated) {\textbf{Tier 1 — Curated memory}\\\scriptsize \texttt{*.md} (canonical)};
\node[super, below=of curated] (super) {\textbf{Supersession sidecars}\\\scriptsize \texttt{*.md.supersessions.jsonl}\\\scriptsize (append-only)};
\node[digest, below=of super] (digest) {\textbf{Tier 2 — Journal digests}\\\scriptsize transcript-distilled summaries};

% --- Dashed boundary grouping the curated + supersession ---
\node[group, fit=(curated)(super), label={[legend,anchor=south west,xshift=0pt,yshift=-2pt]north west:Operator-authored corpus}] (curatedgroup) {};

% --- Merge engine ---
\node[engine, right=2.2cm of super.east, anchor=west] (engine) {%
  \textbf{Retrieval Engine}\\[2pt]
  \scriptsize SQLite FTS5 + BM25\\
  \scriptsize \textsl{accretive merge}\\
  \scriptsize precedence: super $\prec$ canonical
};

% --- LLM consumer ---
\node[consumer, right=1.4cm of engine.east, anchor=west] (llm) {%
  \textbf{LLM}\\[1pt]\scriptsize agent / CDSS
};

% --- File-system watcher (small, top-right of engine) ---
\node[draw, thick, fill=yellow!18, draw=yellow!60!black,
      rectangle, rounded corners=2pt, minimum width=2cm, minimum height=0.7cm,
      above=0.5cm of engine.north, anchor=south,
      font=\sffamily\scriptsize] (watcher) {fs-watcher (1\,s)};

% --- Arrows ---
\draw[flow] (curated.east) -- ++(0.5cm,0) |- (engine.west |- curated);
\draw[flow] (super.east)   -- (engine.west |- super);
\draw[flow] (digest.east)  -- ++(0.5cm,0) |- (engine.west |- digest);

\draw[flow] (engine.east) -- node[above,font=\sffamily\scriptsize]{merged result\\(super first)} (llm.west);

\draw[flow, dashed, gray] (watcher.south) -- (engine.north);
\draw[flow, dashed, gray] (watcher.west) -| ++(-1cm,-0.3cm) -| (curated.north);

% --- Legend ---
\node[legend, below=0.4cm of digest.south, anchor=north, align=left, text width=12cm] {%
  \textbf{Precedence invariant:} for any retrieved file with $k \geq 1$ supersession records,
  the merged output places all $k$ records before the canonical body —
  ensuring sequential readers (LLMs or humans) encounter the correction first.\\
  \textit{Theorem~\ref{thm:precedence}: empirically verified at $100\%$ over $N{=}1000$ randomized queries.}
};

\end{tikzpicture}%
}
\caption{Architecture of memory-oracle. The retrieval engine indexes three corpus
tiers (curated files, supersession sidecars, journal digests) and merges
supersession records ahead of canonical bodies at read time. A filesystem
watcher absorbs new writes within $\sim$1\,s, enabling a self-extending
corpus. The LLM consumer (clinical decision support, coding agent, trading
strategy assessor) receives a single merged document in which corrections
structurally precede the assertions they supersede.}
\label{fig:architecture}
\end{figure}
